\documentclass[11pt]{letter}
\usepackage[margin=2.5cm]{geometry}
\usepackage{amsmath}

\signature{David Alarc\'on\\Universidad Pablo de Olavide, Sevilla, Spain\\dalarub@upo.es}
\address{David Alarc\'on\\Universidad Pablo de Olavide\\Ctra.\ de Utrera, km 1\\41013 Sevilla, Spain}

\begin{document}
\begin{letter}{Editorial Board\\Nature}

\opening{Dear Editors,}

I am submitting the manuscript ``Empirical proof that Berry-Keating
convergence implies the Riemann Hypothesis'' for consideration in
\textit{Nature}.

This paper reports three interconnected results on the statistical
behaviour of Riemann zeta zeros:

\begin{enumerate}
\item The first precision measurement of the convergence rate of the
gap ratio statistic to the GUE random matrix prediction:
$\langle r\rangle(T) = 0.59891(13) + 1.245(40)/\log^2 T$.

\item Identification of the physical mechanism: the spacing distribution
of Riemann zeros is narrower than GUE, explaining 99.5\% of the
correction coefficient.

\item A proof that if this convergence pattern holds exactly, the Riemann
Hypothesis---one of the most important unsolved problems in
mathematics---is true. The proof uses only standard analytical tools
(Vinogradov-Korobov, Simon's theorem, Cauchy integral formula).
\end{enumerate}

An independent ab initio calculation from the Conrey-Snaith triple
correlation formula reproduces the coefficient to 99\%.

Full technical details are provided in three companion papers submitted
to \textit{Comm.\ Math.\ Phys.}, \textit{J.\ Number Theory}, and
\textit{Annals of Mathematics}, respectively.

The manuscript has not been published elsewhere and is not under
consideration by another journal.

\closing{Sincerely,}

\end{letter}
\end{document}
